\documentclass{article}
\usepackage[T1]{fontenc}
\usepackage{lmodern}
\usepackage{graphicx}
\usepackage{amsmath,amssymb}
\usepackage[a4paper,margin=2cm]{geometry}
\usepackage[absolute,overlay]{textpos}
\setlength{\TPHorizModule}{1mm}
\setlength{\TPVertModule}{1mm}
\usepackage[indonesian]{babel}
\usepackage{hyperref}
\setlength{\emergencystretch}{2em}
% unit: Halaman 53

\begin{document}

% === Halaman 53 (llm) ===

Langkah-langkah metode Kasiski:

\begin{enumerate}
  \item Temukan semua kriptogram yang berulang di dalam cipherteks (pesan yang panjang biasanya mengandung kriptogram yang berulang).
  \item Hitung jarak antara kriptogram yang berulang
  \item Hitung semua faktor (pembagi) dari jarak tersebut (faktor pembagi menyatakan panjang kunci yang mungkin ).
  \item Tentukan irisan dari himpunan faktor pembagi tersebut. Nilai yang muncul di dalam irisan menyatakan angka yang muncul pada semua faktor pembagi dari jarak-jarak tersebut . Nilai tersebut mungkin adalah panjang kunci.
\end{enumerate}

\end{document}
